\section{Results}\label{sec:results}

Every Sharpe ratio is reported with its annualised return and maximum drawdown; every number in a table is generated from the stored result files (Appendix~\ref{app:repro}).

\subsection{The information is real, small and replicates}\label{sec:res_info}

Table~\ref{tab:ic_windows} gives the fixed booster's rank IC by window and horizon: weekly $\numICWeeklyDev$ ($t=\numICWeeklyDevT$) on development and $\numICWeeklyChk$ ($t=\numICWeeklyChkT$) on the check window; monthly $\numICMonthlyDev$ / $\numICMonthlyChk$; quarterly $\numICQuarterlyDev$ / $\numICQuarterlyChk$ ($t$ per formation). The first read of 2019--2022, taken before any portfolio cell existed, gave $\numICEsixBooster$ ($t=\numICEsixBoosterT$, 47 months) with ridge at $\numICEsixRidge$; the single hold-out read of 2023 (eleven months, rule frozen in advance) gave $\numICTwentyThree$ ($t=\numICTwentyThreeT$) with ridge at $\numICTwentyThreeRidge$. The sign replicated twice outside the selection data and the linear control replicated with it: the replication is about the characteristic set, not the booster. Three qualifications. Out-of-sample $R^2$ is within $\pm0.001$ of zero on both windows and $+0.012$ on 2023: the signal is ordinal. Composites with no fitted magnitudes (literature signs, textbook five, fitted signs) read IC $+0.020$, $+0.016$ and $+0.025$ over 2009--2022 against the booster's $+0.032$, paired $t\le1.4$: whether fitting adds anything cannot be told on this panel. With per-refit selection on its validation slice the same model class reads $+0.014$, no better than ridge ($\Delta\IC=+0.001$). Figure~\ref{fig:ic_by_year}: the yearly IC is positive in nine of ten development years (2013 is zero to three decimals) and higher on 2019--2023, with 2023 just below 2021 ($\numICYearTwentyOne$). The rise is mostly composition: within the cohort of firms present in 2009 there is no detectable trend after 2014, and the 2019--2022 excess sits in post-2009 entrants (cohort $+0.023$, entrants $+0.070$), whose survivorship filtering differs.

\begin{table}[t]
\centering\footnotesize
\setlength{\tabcolsep}{5pt}
\caption{Rank IC of the fixed booster by window and horizon, with the linear control where run. $t=\overline{\IC}/s_{\IC}\sqrt n$ over the $n$ formation periods, treating periods as independent.}
\label{tab:ic_windows}
\input{tables/tab_ic_windows}
\end{table}

\begin{figure}[t]
\centering
\includegraphics[width=0.66\textwidth]{fig_ic_by_year}
\caption{Yearly mean of the monthly rank IC of the fixed booster used for the out-of-window and hold-out reads, 2009--2023; development, check and hold-out windows shaded, with the mean of each window's yearly values.}
\label{fig:ic_by_year}
\end{figure}

\subsection{No construction turns it into the target Sharpe ratio}\label{sec:res_books}

\paragraph{Grid.}
Table~\ref{tab:p12c_cells} lists the best development cells at their declared costs with their check-window values. The pre-declared selection, the W+M blend on the 500 most liquid names at 2~bps, nets $\numSelDevNet$ on development ($\numSelDevAnn\pct$ per year, $\MaxDD$ $\numSelDevDD\pct$, gross $\numSelDevGross$) and $\numSelChkNet$ on the check window ($\numSelChkAnn\pct$, $\numSelChkDD\pct$, gross $\numSelChkGross$). The pass bar of 1.5 fails; no cell reaches 0.5 on both windows ($\numCellsHalfBoth$ of 55). Figure~\ref{fig:cost_frontier} places every cell by gross and net Sharpe ratio: the cloud sits below 1.5 on the gross axis in both windows, and the best development cell is not the best check cell (cross-window rank correlation of the 55 cells $\numXwinRhoNet$ net, $\numXwinRhoGross$ gross), as selection on a noisy but not empty ordering produces. Table~\ref{tab:precision} gives the precision: with iid standard errors $\numSelDevSE$ and $\numSelChkSE$ the target is $\numSelDevZ$ and $\numSelChkZ$ standard errors away; pooled over 2009--2022 the selected cell nets $\numPooledNet\pm\numPooledSE$ with upper 95\% limit $\numPooledUpper$; none of the three net intervals excludes zero.

\begin{table}[t]
\centering\footnotesize
\setlength{\tabcolsep}{3.5pt}
\caption{Best development cells of the 55-cell grid at the declared cost of their universe (ALL 15, LIQ 5, TOP500 2 bps per side), development $\MaxDD\ge-25\%$. $\dagger$ marks the pre-declared selection. Last row: counts over all 55 cells.}
\label{tab:p12c_cells}
\fittab{tables/tab_p12c_cells}
\end{table}

\begin{figure}[t]
\centering
\includegraphics[width=0.92\textwidth]{fig_cost_frontier}
\caption{Gross against net Sharpe ratio at the declared cost for the 55 cells, development (left) and check window (right); colour by universe, marker by aim family, the selected cell outlined; dashed lines at 1.5, dotted line net = gross.}
\label{fig:cost_frontier}
\end{figure}

\begin{table}[t]
\centering\footnotesize
\setlength{\tabcolsep}{4pt}
\caption{Precision of the selected cell's Sharpe ratios: iid standard error (eq.~\ref{eq:se}), 95\% interval, distance to the target 1.5 and to zero in standard errors; counts over the 55 cells of pooled upper limits reaching 1.0 and 1.5.}
\label{tab:precision}
\input{tables/tab_precision}
\end{table}

\paragraph{Selection bar and transfer.}
The expected maximum of $N$ independent zero-mean null Sharpe ratios over ten years is $\numExpMaxTwoNineFour$ at $N=294$ (stored) and $\numExpMaxThreeFourNine$ at $N=349$ (simulated); the best development net Sharpe ratio on the rebuilt panels, $\numSelDevNet$ (0.83 on the superseded older build), is below it, so the net books alone cannot be told from selection noise, while the gross books (selected $\numSelDevGross$, best $\numMaxGrossDev$) and the weekly IC's $t$ of $\numICWeeklyDevT$ are above it: the distance between the information and the net book is partly transfer and partly cost. The bar is indicative: the cells are correlated, which lowers it, and a no-information net Sharpe ratio has mean $-\mathrm{drag}/\sigma$. The annualised ICIR of the raw weekly prediction ($t\sqrt{52/n}$) is $\numWeeklyRankICIRRefitDev$ on development and $\numWeeklyRankICIRRefitChk$ on the check window; the highest gross Sharpe ratio any of the 55 constructions realises is $\numMaxGrossDev$ and $\numMaxGrossChk$, a measured transfer coefficient of $\numTransferCoefRefitDev$ and $\numTransferCoefRefitChk$. Both are upper estimates (the numerator is a maximum over 55 cells), the check ratio carries $\pm\numTransferCoefChkSE$ from the Sharpe standard error alone, and the book is volatility-scaled and neutralised over the tradable set while the ICIR scores every name, so the ratio is an ex-post analogue of a transfer coefficient. Gross rises with trading and net falls with it: the weekly-traded $\theta=1$ books (R1, R4, R4b) reach the highest gross Sharpe ratios at one to two turns a week; on development they are net-positive at 1--2 bps, between $\numThetaOneFiveBpsMin$ and $\numThetaOneFiveBpsMax$ at 5 bps and deeply negative at 15 bps; the slow ALL books ($\theta=0.05$) keep half to two thirds of the development gross at under a tenth of the turnover (Table~\ref{tab:p12c_cells}, Figure~\ref{fig:cost_frontier}).

\paragraph{Crisis years in the check window.}
The check window contains the March 2020 crash and the 2022 bear market. Table~\ref{tab:crisis} splits it by calendar year for three series whose yearly values are stored: the monthly rank IC, the monthly decile book of Section~\ref{sec:power}, and a full-breadth rank-weighted book made dollar-, beta- and SIC-division-neutral and held in three overlapping monthly cohorts (both ALL, 15 bps); the weekly grid books are stored only as window summaries. The 2020 yearly IC, $\numICYearTwenty$, is close to the development mean $\numICDevYearMean$ with standard error $\numICTwentySE$: the year neither confirms nor excludes a failure. The crash months are negative (February $\numICFebTwenty$, the $\numICFebTwentyRank$th lowest of 168 months; March $\numICMarTwenty$) and April recovers ($\numICAprTwenty$); the decile book lost $\numDecileFebTwenty\pct$ and $\numDecileMarTwenty\pct$, gained $\numDecileAprTwenty\pct$ in April and ended 2020 at $\numDecileYearTwenty\pct$. The weakest check year for the IC is 2022 ($\numICYearTwentyTwo$); without 2020 the check mean of yearly ICs moves from $\numICChkYearMean$ to $\numICChkYearMeanExTwenty$. The decile book's worst year is $\numDecileWorstYear$ ($\numDecileWorstYearNet\pct$), inside the development window; the neutral book's worst year is 2020 ($\numNeutralYearTwenty\pct$, against $\numNeutralYearTwentyTwo\pct$ in 2022), so the size and sign of a crisis-year loss depend on the construction. The check window was the better window for the information and for the decile book's return ($\numDecileChkYearMean\pct$ per year against $\numDecileDevYearMean\pct$ on development), not for its risk: volatility rises from $\numDecileVolDev\pct$ to $\numDecileVolChk\pct$. The weekly grid summaries also point to a higher-volatility rather than a lower-return window: for the selected cell gross return is $\numSelDevGrossAnn\pct$ per year on development and $\numSelChkGrossAnn\pct$ on the check window while volatility rises from $\numSelDevVol\pct$ to $\numSelChkVol\pct$; across the 55 cells the check-to-development volatility ratio is $\numVolRatioMin$--$\numVolRatioMax$ (median $\numVolRatioMedian$), gross return is higher on the check window in $\numCellsHigherChkGrossAnn$ cells, and $\numCellsDeeperChkDD$ cells have a deeper check-window drawdown at the declared cost despite the shorter window. This is consistent with a crisis effect on the weekly Sharpe ratios but does not establish one. The ex-2020 figures are a post-hoc diagnostic and the pass rule is unchanged; a predictor restricted to ordinary conditions would need an ex-ante trigger (volatility or drawdown), a new pre-declared trial that was not declared.

\begin{table}[t]
\centering\footnotesize
\setlength{\tabcolsep}{5pt}
\caption{The check window by calendar year (post-hoc split): yearly mean of the monthly rank IC of the fixed booster, and calendar-year net return of the monthly decile long--short book (ALL, 15 bps) and of a neutral rank-weighted book held in three-month cohorts.}
\label{tab:crisis}
\input{tables/tab_crisis}
\end{table}

\subsection{What a net Sharpe ratio of 1.5 would take}\label{sec:res_takes}

Equation~\eqref{eq:ladder} reproduces the stored cost ladder of all 55 cells to within 0.02, so its inversion, eq.~\eqref{eq:ic_mult}, is arithmetic on the stored summaries (Table~\ref{tab:what_it_takes}). No cell's gross Sharpe ratio reaches 1.5 on either window (maximum $\numMaxGrossDev$ development, $\numMaxGrossChk$ check), so no cost level, including zero, closes the gap; the cost that would is negative for every cell. The information multiple is $\numKappaSelDev$ (development) and $\numKappaSelChk$ (check) for the selected cell and $\numKappaLeastBoth$ for the least demanding cell on both windows: the shortfall is information, not cost. It is a floor, not a recipe: between the two windows the least demanding cell's weekly IC doubled ($\numICWeeklyDev$ to $\numICWeeklyChk$) while its gross Sharpe ratio stayed at $\numLeastBothDevGross$ and $\numLeastBothChkGross$. The selected cell's break-even cost on development is $\numBreakevenSelDev$ bps per side; at its declared 2 bps the drag takes about four tenths of the gross.

\begin{table}[t]
\centering\footnotesize
\setlength{\tabcolsep}{3.5pt}
\caption{Cost-ladder inversion (eq.~\ref{eq:ic_mult}) for the selected cell and the least demanding cells: annual volatility $\sigma$, turns per year, declared cost $c_d$ and break-even cost $\hat c$ per side, the gross Sharpe ratio $S_g^{need}$ needed for net 1.5 at $c_d$, and the information multiple $\kappa$.}
\label{tab:what_it_takes}
\fittab{tables/tab_what_it_takes}
\end{table}

\subsection{Graph networks trained end to end add no information}\label{sec:res_graphs}

\paragraph{Levels.}
Table~\ref{tab:p13_ic} gives the seed-mean rank IC of each graph cell by horizon and window, Table~\ref{tab:p13_books} the selected book of each prediction set under the identical grid and rules. The no-graph baseline C0 nets $\numPthirteenCzeroDevNet$ / $\numPthirteenCzeroChkNet$ (development / check); C1 $\numPthirteenConeDevNet$ / $\numPthirteenConeChkNet$, C2 $\numPthirteenCtwoDevNet$ / $\numPthirteenCtwoChkNet$, C3 $\numPthirteenCthreeDevNet$ / $\numPthirteenCthreeChkNet$; the booster on the same grid $\numSelDevNet$ / $\numSelChkNet$. On development every graph book is at or below the baseline (on the check window C3 is above it, $\numPthirteenCthreeChkNet$ against $\numPthirteenCzeroChkNet$, inside one standard error of $\numSelChkSE$); all are far from 1.5.

\paragraph{Paired $\Delta\IC$.}
Table~\ref{tab:p13_dic} and Figure~\ref{fig:p13_dic} are the primary evidence. On development the demeaned $\Delta\IC$ of C1 against C0 is small and of mixed sign across horizons; C2 and C3 are slightly negative on the weekly horizon; the largest $|t|$ against C0 on development is $\numPthirteenMaxAbsDicT$, below the pre-declared threshold of 2. The only comparison with $t\ge2$ (C1 against the aggregate-feature control C0g, weekly; the only other $|t|\ge2$ is negative, C1 against C0 on the quarterly check window, raw basis) is not a pass condition and does not repeat on the check window.

\paragraph{Rewired null and label shift.}
C4, where it ran, has the highest monthly development IC of any graph cell and quarterly and check-window ICs inside the range of the real-graph cells, as expected if message passing fits degree and industry composition rather than specific pairings; its paired comparison was not completed. None of the four label-shift diagnostics meets the persistence rule, but at these ICs the diagnostic has little power (the real development ICs have $t\le1.7$; on the quarterly horizon the shifted models read higher ICs than the real ones): it neither attributes nor clears a signal that is not there.

\paragraph{Scope.}
Not completed before the pause: the weekly C4 and C5 predictions, the C4 and C5 books, every paired $\Delta\IC$ against C4, all of C5's paired $\Delta\IC$, the weekly label-shift runs and the formal selection script. Because a pass needs a development $\Delta\IC$ against C0 with $t\ge2$ on every horizon used, and no horizon of C1--C3 has one, the missing C4 comparisons cannot make C1--C3 pass. C5 has neither a book nor a paired $\Delta\IC$; its monthly and quarterly ICs sit at the baseline, so a pass is unlikely but not excluded. For the completed cells the stop rule's negative form applies: graph models trained end to end add a $\Delta\IC$ indistinguishable from zero over the same features without a graph, and the gap to the target remains the information gap of Section~\ref{sec:res_takes}.

\begin{table}[t]
\centering\footnotesize
\setlength{\tabcolsep}{3pt}
\caption{Graph programme: seed-mean rank IC by cell and horizon, development $(t)$ / check $(t)$; naive per-formation $t$. ``pending'' marks cells not run before the pause.}
\label{tab:p13_ic}
\input{tables/tab_p13_ic}
\end{table}

\begin{table}[t]
\centering\footnotesize
\setlength{\tabcolsep}{3pt}
\caption{Graph programme: the selected book of each prediction set under the 55-cell grid and rules of Section~\ref{sec:books}; net Sharpe ratio (annualised return, $\MaxDD$), gross Sharpe ratio, standard errors (iid / 13-week block bootstrap). Eligible cells are those the rule may select.}
\label{tab:p13_books}
\fittab{tables/tab_p13_books}
\end{table}

\begin{table}[t]
\centering\footnotesize
\setlength{\tabcolsep}{3.5pt}
\caption{Paired $\Delta\IC$ of each graph cell against the controls (C0 no-graph; C0g aggregate features), by horizon and window, demeaned and raw bases, month-level $t$ in parentheses; $*$ marks development demeaned $|t|\ge2$.}
\label{tab:p13_dic}
\fittab{tables/tab_p13_dic}
\end{table}

\begin{figure}[t]
\centering
\includegraphics[width=0.64\textwidth]{fig_p13_dic}
\caption{Development paired $\Delta\IC$ (demeaned basis) of the graph cells against C0 by horizon, with $\pm1.96$ standard errors recovered from the $t$-statistic (approximate: $t$ is stored to one decimal); hollow markers are the check window.}
\label{fig:p13_dic}
\end{figure}

\begin{table}[t]
\centering\footnotesize
\setlength{\tabcolsep}{4pt}
\caption{Label-shift diagnostic (seed 0): IC of the model trained on labels shifted back one year against the real model; the declared persistence rule (shifted development IC $>0$, $t\ge2$, at least half the real IC) is not met in any case.}
\label{tab:label_shift}
\fittab{tables/tab_label_shift}
\end{table}
